\makereporttitle{2024}{septembrie}

\makereportpreamble

Pe parcursul acestei luni, am reluat noțiunile studiate pe zona de haos cuantic, în vederea continuării cercetării în această direcție și pentru a redacta referatul științific în care să descriu progresul și realizările după primul an de studii. Referatul conține o introducere/recapitulare a unor noțiuni de mecanică clasică, de mecanică analitică (formalismul Lagrangean, respectiv cel Hamiltonian), de mecanică cuantică (formalismul bazat pe operatori), sisteme dinamice continue (inclusiv rezolvarea lor numerică), teoria matricilor aleatoare, factorizări de matrici (în vederea determinării nivelelor de energie) și haos cuantic (la nivel de perspectivă).

Luna aceasta am finalizat și transmis proiectul aferent cursului doctoral de „\emph{Metodologia cercetării științifice și prelucrarea statistică a datelor”}, ținut de conf.~dr.~Cristian Necula. Am folosit \emph{wavelet analysis} pentru a analiza un set de date de vreme, înregistrate pe parcursul a 10 ani în orașul Mannheim din Germania. Codul sursă aferent proiectului poate fi găsit \href{https://github.com/GabrielMajeri/wavelet-analysis-project}{pe GitHub}.

Pentru a înțelege codul de calcul numeric pe care l-am primit de la Arnaud Leclerc și Georges Jolicard, autorii lucrării „\emph{Calculating eigenvalues and eigenvectors of parameter-dependent hamiltonians using an adaptative wave operator method}”, am început să lucrez exerciții de MATLAB de pe \href{https://www.mathworks.com/matlabcentral/cody/problems?term=&sort=created+asc}{MATLAB Cody}, platforma de exerciții oferită de \href{https://www.mathworks.com/}{MathWorks}. De asemenea, având în vedere că mare parte din codul de simulare numerică dezvoltat anterior de domnul profesor Băran este scris în \href{https://fortran-lang.org/learn/}{Fortran}, am început să învăț și acest limbaj de programare.

În vederea identificării a noi direcții/subiecte de cercetare, am urmărit video-ul \href{https://www.youtube.com/watch?v=JoFW2uSd3Uo}{\emph{Physics Informed Machine Learning: High Level Overview of AI and ML in Science and Engineering}} a lui Steve Brunton. De aici am aflat despre arhitectura de rețea neurală denumită \href{https://arxiv.org/abs/2003.04630}{Lagrangian Neural Network}, care poate fi folosită pentru a simula evoluția unor sisteme dinamice cu ajutorul învățării automate.

\vspace{0.5em}

\textbf{Observațiile conducătorului științific}:

\vspace{3em}

\makereportpostamble{24 septembrie 2024}
